%%%%%%%%%%%%%%%%%%%%%%%%%%%%%%%%%%%%%%%%%%%%%%%%%%%%%%
% IDIOMA
%%%%%%%%%%%%%%%%%%%%%%%%%%%%%%%%%%%%%%%%%%%%%%%%%%%%%%

% es-tabla       -> "Tabla" en vez de "Cuadro"
% es-nodecimaldot -> el punto decimal se queda como punto (igual que en R)
% es-noquoting   -> evita conflictos de < y > con TikZ
\usepackage[english,spanish,es-tabla,es-nodecimaldot,es-noquoting]{babel}

%%%%%%%%%%%%%%%%%%%%%%%%%%%%%%%%%%%%%%%%%%%%%%%%%%%%%%
% GEOMETRÍA
%%%%%%%%%%%%%%%%%%%%%%%%%%%%%%%%%%%%%%%%%%%%%%%%%%%%%%

\usepackage[
top=0.65in,
bottom=0.65in,
left=0.65in,
right=0.65in
]{geometry}

%%%%%%%%%%%%%%%%%%%%%%%%%%%%%%%%%%%%%%%%%%%%%%%%%%%%%%
% ENCABEZADO Y PIE DE PÁGINA (estilo artículo)
%%%%%%%%%%%%%%%%%%%%%%%%%%%%%%%%%%%%%%%%%%%%%%%%%%%%%%

\usepackage{fancyhdr}
\setlength{\headheight}{15pt}

% Páginas normales: título corto a la izquierda,
% curso a la derecha y línea roja, como en la revista
\pagestyle{fancy}
\fancyhf{}
\fancyhead[L]{\footnotesize\AncizarSansRegular\color{grisOscuroUN}\titulocorto}
\fancyhead[R]{\footnotesize\AncizarSansRegular\color{grisOscuroUN}\materia\ \textbullet\ \semestre}
\fancyfoot[C]{\AncizarSansRegular\thepage}
\renewcommand{\headrulewidth}{0.6pt}
\renewcommand{\headrule}{{\color{rojoUN}\hrule width\headwidth height\headrulewidth}}

% Primera página del artículo: solo el número de página
\fancypagestyle{plain}{%
  \fancyhf{}
  \fancyfoot[C]{\AncizarSansRegular\thepage}
  \renewcommand{\headrulewidth}{0pt}%
}


%%%%%%%%%%%%%%%%%%%%%%%%%%%%%%%%%%%%%%%%%%%%%%%%%%%%%%
% FUENTES
%%%%%%%%%%%%%%%%%%%%%%%%%%%%%%%%%%%%%%%%%%%%%%%%%%%%%%

\usepackage{fontspec}
\usepackage{microtype}

%%%%%%%%%%%%%%%%%%%%%%%%%%%%%%%%%%%%%%%%%%%%%%%%%%%%%%
% MATEMÁTICAS
%%%%%%%%%%%%%%%%%%%%%%%%%%%%%%%%%%%%%%%%%%%%%%%%%%%%%%

\usepackage{amsmath}
\usepackage{amssymb}
\usepackage{amsthm}
\usepackage{cancel}
\usepackage{bm}        % letras griegas en negrita: \bm{\beta}

%%%%%%%%%%%%%%%%%%%%%%%%%%%%%%%%%%%%%%%%%%%%%%%%%%%%%%
% GRÁFICOS Y TABLAS
%%%%%%%%%%%%%%%%%%%%%%%%%%%%%%%%%%%%%%%%%%%%%%%%%%%%%%

\usepackage{graphicx}
\usepackage{booktabs}
\usepackage{tabularx}
\usepackage{multirow}
\usepackage{multicol}
\usepackage{float}
\usepackage{pdflscape}

%%%%%%%%%%%%%%%%%%%%%%%%%%%%%%%%%%%%%%%%%%%%%%%%%%%%%%
% COLORES
%%%%%%%%%%%%%%%%%%%%%%%%%%%%%%%%%%%%%%%%%%%%%%%%%%%%%%

\usepackage[dvipsnames,table]{xcolor}
\usepackage{colortbl}
%==================================================
% COLORES INSTITUCIONALES UNAL
%==================================================

\definecolor{verdeUN}{RGB}{148,180,59}      % Pantone 376C
\definecolor{rojoUN}{RGB}{166,28,49}        % Pantone 187C
\definecolor{rojoOscuroUN}{RGB}{118,35,47}  % Pantone 188C
\definecolor{verdeOscuroUN}{RGB}{70,107,63} % Pantone 7743C
\definecolor{grisOscuroUN}{RGB}{86,90,92}   % Pantone 425C
\definecolor{grisClaroUN}{RGB}{177,178,176} % Pantone 421C

%==================================================
% COLORES PARA CÓDIGO
%==================================================

\definecolor{bgcream}{HTML}{FFF8E7}
\definecolor{keywordcolor}{RGB}{70,107,63}
\definecolor{stringcolor}{RGB}{166,28,49}
\definecolor{commentcolor}{RGB}{86,90,92}

%%%%%%%%%%%%%%%%%%%%%%%%%%%%%%%%%%%%%%%%%%%%%%%%%%%%%%
% TIKZ
%%%%%%%%%%%%%%%%%%%%%%%%%%%%%%%%%%%%%%%%%%%%%%%%%%%%%%

\usepackage{tikz}

%%%%%%%%%%%%%%%%%%%%%%%%%%%%%%%%%%%%%%%%%%%%%%%%%%%%%%
% HIPERVÍNCULOS
%%%%%%%%%%%%%%%%%%%%%%%%%%%%%%%%%%%%%%%%%%%%%%%%%%%%%%

\usepackage{xurl}
\usepackage[hidelinks]{hyperref}

%%%%%%%%%%%%%%%%%%%%%%%%%%%%%%%%%%%%%%%%%%%%%%%%%%%%%%
% REFERENCIAS (igual que la plantilla del profesor)
%%%%%%%%%%%%%%%%%%%%%%%%%%%%%%%%%%%%%%%%%%%%%%%%%%%%%%
% \citet{clave}  -> Autor (año)
% \citep{clave}  -> (Autor, año)

%\usepackage[round]{natbib}
%\bibliographystyle{abbrvnat-es}   % abbrvnat en español (archivo abbrvnat-es.bst)
%\setlength{\bibsep}{4pt}

%%%%%%%%%%%%%%%%%%%%%%%%%%%%%%%%%%%%%%%%%%%%%%%%%%%%%%
% CAJAS
%%%%%%%%%%%%%%%%%%%%%%%%%%%%%%%%%%%%%%%%%%%%%%%%%%%%%%

\usepackage[most]{tcolorbox}

%%%%%%%%%%%%%%%%%%%%%%%%%%%%%%%%%%%%%%%%%%%%%%%%%%%%%%
% CAPTIONS Y LISTAS
%%%%%%%%%%%%%%%%%%%%%%%%%%%%%%%%%%%%%%%%%%%%%%%%%%%%%%

\usepackage{caption}
\usepackage{enumitem}

\captionsetup{font=small, labelfont=bf, labelsep=period, margin=0.5cm}
\setlength{\belowcaptionskip}{6pt}

%%%%%%%%%%%%%%%%%%%%%%%%%%%%%%%%%%%%%%%%%%%%%%%%%%%%%%
% NÚMEROS
%%%%%%%%%%%%%%%%%%%%%%%%%%%%%%%%%%%%%%%%%%%%%%%%%%%%%%

\usepackage{siunitx}

\sisetup{
table-number-alignment = center,
input-symbols = {+},
table-space-text-post = {+}
}

%%%%%%%%%%%%%%%%%%%%%%%%%%%%%%%%%%%%%%%%%%%%%%%%%%%%%%
% SECCIONES
%%%%%%%%%%%%%%%%%%%%%%%%%%%%%%%%%%%%%%%%%%%%%%%%%%%%%%

\usepackage{titlesec}

\titleformat{\section}
{\Large\bfseries}
{\thesection.}{0.5em}{}
[\vspace{0.3em}\color{rojoUN}\titlerule]

\titleformat{\subsection}
{\large\bfseries\color{verdeOscuroUN}}
{\thesubsection.}{0.5em}{}

\titleformat{\subsubsection}
{\normalsize\bfseries\itshape}
{\thesubsubsection.}{0.5em}{}

\titlespacing*{\section}{0pt}{2.2ex plus 0.5ex}{1.2ex}
\titlespacing*{\subsection}{0pt}{1.8ex plus 0.5ex}{0.8ex}

%%%%%%%%%%%%%%%%%%%%%%%%%%%%%%%%%%%%%%%%%%%%%%%%%%%%%%
% COLORBOX
%%%%%%%%%%%%%%%%%%%%%%%%%%%%%%%%%%%%%%%%%%%%%%%%%%%%%%

\newtcolorbox{cuadro_verde}[1]{
    enhanced,
    breakable,
    colback=white,
    colframe=verdeUN,
    colbacktitle=verdeUN,
    coltitle=white,
    fonttitle=\AncizarSansSemiBold,
    title=#1,
    boxrule=0.8pt,
    arc=2mm,
    left=3mm,
    right=3mm,
    top=2mm,
    bottom=2mm
}

\newtcolorbox{cuadro_verdeOscuro}[1]{
    enhanced,
    breakable,
    colback=white,
    colframe=verdeOscuroUN,
    colbacktitle=verdeOscuroUN,
    coltitle=white,
    fonttitle=\AncizarSansSemiBold,
    title=#1,
    boxrule=0.8pt,
    arc=2mm,
    left=3mm,
    right=3mm,
    top=2mm,
    bottom=2mm
}

\newtcolorbox{cuadro_rojo}[1]{
    enhanced,
    breakable,
    colback=white,
    colframe=rojoUN,
    colbacktitle=rojoUN,
    coltitle=white,
    fonttitle=\AncizarSansSemiBold,
    title=#1,
    boxrule=0.8pt,
    arc=2mm,
    left=3mm,
    right=3mm,
    top=2mm,
    bottom=2mm
}

\newtcolorbox{cuadro_rojo_Oscuro}[1]{
    enhanced,
    breakable,
    colback=white,
    colframe=rojoOscuroUN,
    colbacktitle=rojoOscuroUN,
    coltitle=white,
    fonttitle=\AncizarSansSemiBold,
    title=#1,
    boxrule=0.8pt,
    arc=2mm,
    left=3mm,
    right=3mm,
    top=2mm,
    bottom=2mm
}

\newtcolorbox{cuadro_grisClaro}[1]{
    enhanced,
    breakable,
    colback=white,
    colframe=grisClaroUN,
    colbacktitle=grisClaroUN,
    coltitle=white,
    fonttitle=\AncizarSansSemiBold,
    title=#1,
    boxrule=0.8pt,
    arc=2mm,
    left=3mm,
    right=3mm,
    top=2mm,
    bottom=2mm
}

\newtcolorbox{cuadro_grisOscuro}[1]{
    enhanced,
    breakable,
    colback=white,
    colframe=grisOscuroUN,
    colbacktitle=grisOscuroUN,
    coltitle=white,
    fonttitle=\AncizarSansSemiBold,
    title=#1,
    boxrule=0.8pt,
    arc=2mm,
    left=3mm,
    right=3mm,
    top=2mm,
    bottom=2mm
}


\newtcolorbox{nota}{
    enhanced,
    breakable,
    colback=white,
    colframe=verdeUN,
    borderline west={3pt}{0pt}{verdeUN},
    boxrule=0pt,
    arc=0mm,
    left=5mm,
    right=3mm,
    top=2mm,
    bottom=2mm
}


%%%%%%%%%%%%%%%%%%%%%%%%%%%%%%%%%%%%%%%%%%%%%%%%%%%%%%
% LISTINGS
%%%%%%%%%%%%%%%%%%%%%%%%%%%%%%%%%%%%%%%%%%%%%%%%%%%%%%

%%%%%%%%%%%%%%%%%%%%%%%%%%%%%%%%%%%%%%%%%%%%%%%%%%%%%%
% FUENTE JETBRAINS MONO
%%%%%%%%%%%%%%%%%%%%%%%%%%%%%%%%%%%%%%%%%%%%%%%%%%%%%%

% Usa IBM Plex Mono si está instalada; si no (por ejemplo
% en otro computador), usa la JetBrains Mono que viene en
% Config/Fuentes, para que el documento siempre compile.
\IfFontExistsTF{IBM Plex Mono}{%
  \setmonofont{IBM Plex Mono}%
}{%
  \setmonofont{JetBrainsMono}[
  Path=Config/Fuentes/JetBrains/,
  Extension=.ttf,
  UprightFont=*-Regular,
  ItalicFont=*-Italic,
  BoldFont=JetBrainsMonoNL-Bold,
  BoldItalicFont=JetBrainsMonoNL-BoldItalic,
  Scale=0.9
  ]%
}


\usepackage{listings}

\lstdefinestyle{Rcode}{
language=R,
backgroundcolor=\color{bgcream},
keywordstyle=\color{keywordcolor}\bfseries,
commentstyle=\color{commentcolor}\itshape,
stringstyle=\color{stringcolor},
basicstyle=\fontsize{10}{14}\ttfamily\selectfont,
numbers=left,
numberstyle=\footnotesize\color{grisClaroUN},
numbersep=8pt,
frame=single,
rulecolor=\color{grisClaroUN},
breaklines=true,
showstringspaces=false,
morekeywords={library,ggplot,df2,zcounts},
xleftmargin=10pt,
framesep=6pt
}

\lstdefinestyle{Routput}{
basicstyle=\fontsize{10}{14}\ttfamily\selectfont,
backgroundcolor=\color{white},
frame=single,
rulecolor=\color{grisClaroUN},
breaklines=true,
showstringspaces=false,
numbers=none,
xleftmargin=10pt,
framesep=6pt
}

%%%%%%%%%%%%%%%%%%%%%%%%%%%%%%%%%%%%%%%%%%%%%%%%%%%%%%
% FUENTE PRINCIPAL DEL DOCUMENTO (SERIF)
%%%%%%%%%%%%%%%%%%%%%%%%%%%%%%%%%%%%%%%%%%%%%%%%%%%%%%

% (sintaxis estándar de fontspec: funciona también con
% versiones de TeX Live más viejas)
\setmainfont{AncizarSerif}[
Path=Config/Fuentes/Serif/,
Extension=.ttf,
UprightFont=*-Regular,
ItalicFont=*-Italic,
BoldFont=*-Bold,
BoldItalicFont=*-BoldItalic
]

%%%%%%%%%%%%%%%%%%%%%%%%%%%%%%%%%%%%%%%%%%%%%%%%%%%%%%
% FUENTES SANS (PORTADA)
%%%%%%%%%%%%%%%%%%%%%%%%%%%%%%%%%%%%%%%%%%%%%%%%%%%%%%

\newfontfamily\AncizarSansMedium[
Path=Config/Fuentes/Sans/
]
{AncizarSans-Medium.ttf}

\newfontfamily\AncizarSansLight[
Path=Config/Fuentes/Sans/
]
{AncizarSans-Light.ttf}

\newfontfamily\AncizarSansLightItalic[
Path=Config/Fuentes/Sans/
]
{AncizarSans-LightItalic.ttf}

\newfontfamily\AncizarSansRegular[
Path=Config/Fuentes/Sans/
]
{AncizarSans-Regular.ttf}

\newfontfamily\AncizarSansSemiBold[
Path=Config/Fuentes/Sans/
]
{AncizarSans-SemiBold.ttf}

\newfontfamily\AncizarSansExtraBold[
Path=Config/Fuentes/Sans/
]
{AncizarSans-ExtraBold.ttf}

%%%%%%%%%%%%%%%%%%%%%%%%%%%%%%%%%%%%%%%%%%%%%%%%%%%%%%
% FUENTES SERIF AUXILIARES
%%%%%%%%%%%%%%%%%%%%%%%%%%%%%%%%%%%%%%%%%%%%%%%%%%%%%%

\newfontfamily\AncizarSerifMedium[
Path=Config/Fuentes/Serif/
]
{AncizarSerif-Medium.ttf}

\newfontfamily\AncizarSerifLight[
Path=Config/Fuentes/Serif/
]
{AncizarSerif-Light.ttf}

\newfontfamily\AncizarSerifLightItalic[
Path=Config/Fuentes/Serif/
]
{AncizarSerif-LightItalic.ttf}

\newfontfamily\AncizarSerifRegular[
Path=Config/Fuentes/Serif/
]
{AncizarSerif-Regular.ttf}

\newfontfamily\AncizarSerifSemiBold[
Path=Config/Fuentes/Serif/
]
{AncizarSerif-SemiBold.ttf}

\newfontfamily\AncizarSerifExtraBold[
Path=Config/Fuentes/Serif/
]
{AncizarSerif-ExtraBold.ttf}

%%%%%%%%%%%%%%%%%%%%%%%%%%%%%%%%%%%%%%%%%%%%%%%%%%%%%%
% RESUMEN Y ABSTRACT (estructura de la plantilla
% del profesor, con el estilo de cajas de la mía)
%%%%%%%%%%%%%%%%%%%%%%%%%%%%%%%%%%%%%%%%%%%%%%%%%%%%%%
% Uso:
% \begin{resumen}
%   texto...
%   \palabrasclave{palabra 1, palabra 2}
% \end{resumen}

\newcommand{\palabrasclave}[1]{%
  \par\medskip\noindent{\AncizarSansSemiBold\color{rojoUN}Palabras clave:}\ #1.}
\newcommand{\keywords}[1]{%
  \par\medskip\noindent{\AncizarSansSemiBold\color{rojoUN}Keywords:}\ #1.}

\newtcolorbox{cajaresumen}[1]{
    enhanced,
    breakable,
    colback=white,
    colframe=verdeUN,
    borderline west={3pt}{0pt}{verdeUN},
    boxrule=0pt,
    arc=0mm,
    left=5mm, right=3mm, top=2mm, bottom=2mm,
    fontupper=\small,
    title={#1},
    coltitle=rojoUN,
    colbacktitle=white,
    fonttitle=\AncizarSansSemiBold\large,
    attach title to upper={\par\smallskip}
}

\newenvironment{resumen}{\begin{cajaresumen}{Resumen}}{\end{cajaresumen}}
\renewenvironment{abstract}{\begin{cajaresumen}{Abstract}\selectlanguage{english}}{\end{cajaresumen}}

%%%%%%%%%%%%%%%%%%%%%%%%%%%%%%%%%%%%%%%%%%%%%%%%%%%%%%
% MODO BORRADOR: notas de lo que falta
%%%%%%%%%%%%%%%%%%%%%%%%%%%%%%%%%%%%%%%%%%%%%%%%%%%%%%
% \pendiente{...} muestra una caja roja mientras
% \borradortrue esté activo en main.tex. Con
% \borradorfalse desaparecen todas de una vez.

\newif\ifborrador
\borradortrue

\newcommand{\pendiente}[1]{%
  \ifborrador
  \begin{cuadro_rojo}{Pendiente}#1\end{cuadro_rojo}%
  \fi}

% Enlace al repositorio (el texto viene de Config/datos.tex)
\newcommand{\enlacerepositorio}{\expandafter\url\expandafter{\repositorio}}

% Parámetros de las figuras flotantes: permiten poner figuras grandes
% junto con texto en la misma página (menos páginas medio vacías)
\renewcommand{\topfraction}{0.9}
\renewcommand{\bottomfraction}{0.8}
\renewcommand{\textfraction}{0.07}
\renewcommand{\floatpagefraction}{0.75}
\setcounter{topnumber}{3}
\setcounter{totalnumber}{4}

% Las figuras no se salen de su sección (no terminan después de las referencias)
\usepackage[section]{placeins}
